% ============================================================================
% ARBEITSBLATT DS 4 - Gruppe A: Punto final + Evaluación (S. 33-34)
% ============================================================================

\documentclass[11pt,a4paper]{article}

\usepackage[utf8]{inputenc}
\usepackage[T1]{fontenc}
\usepackage[ngerman]{babel}
\usepackage[left=2cm,right=2cm,top=1.8cm,bottom=1.5cm]{geometry}
\usepackage{helvet}
\usepackage[table]{xcolor}
\usepackage{tabularx}
\usepackage{array}
\usepackage{enumitem}
\usepackage{tcolorbox}
\usepackage{fancyhdr}
\usepackage{lastpage}

\renewcommand{\familydefault}{\sfdefault}

\definecolor{spanisch}{HTML}{c41e3a}
\definecolor{grau}{HTML}{666666}
\definecolor{hellgrau}{HTML}{f5f5f5}

\pagestyle{fancy}
\fancyhf{}
\fancyhead[L]{\small\textcolor{grau}{Spanisch Spätbeginner -- Gruppe A}}
\fancyhead[R]{\small\textcolor{grau}{AB-04-A}}
\fancyfoot[C]{\small\thepage/\pageref{LastPage}}
\renewcommand{\headrulewidth}{0.5pt}

\newcommand{\luecke}[1]{\underline{\hspace{#1}}}
\newcommand{\punktebox}[1]{\hfill\fbox{\small \_\_/#1}}
\newcommand{\es}[1]{\textit{#1}}

\newtcolorbox{merkkasten}[1][]{
    colback=white,
    colframe=spanisch,
    fonttitle=\bfseries\color{spanisch},
    title=#1,
    boxrule=1.5pt,
    arc=0pt,
    left=5pt, right=5pt, top=5pt, bottom=5pt
}

\newtcolorbox{buchverweis}{
    colback=hellgrau,
    colframe=grau,
    boxrule=0.5pt,
    arc=0pt,
    left=5pt, right=5pt, top=3pt, bottom=3pt
}

\newtcolorbox{hilfskasten}{
    colback=hellgrau,
    colframe=spanisch!50,
    boxrule=0.5pt,
    arc=0pt,
    left=5pt, right=5pt, top=3pt, bottom=3pt
}

\newcounter{aufgabe}
\newcommand{\aufgabe}[2]{%
    \stepcounter{aufgabe}%
    \vspace{0.8em}%
    \noindent\textbf{\theaufgabe.} #1 \punktebox{#2}%
    \vspace{0.3em}%
}

\begin{document}

% ============================================================================
% KOPFBEREICH
% ============================================================================
\noindent
\begin{tabularx}{\textwidth}{@{}Xr@{}}
    {\Large\textbf{\textcolor{spanisch}{Punto final + Evaluación}}} & Name: \luecke{5cm} \\[0.3em]
    {\small DS 4 -- Abschluss Unidad 1-2 (S. 33--34)} & Datum: \luecke{3cm} \\
\end{tabularx}

\vspace{0.3em}
\hrule
\vspace{0.8em}

% ============================================================================
% MERKKASTEN: E-Mail schreiben
% ============================================================================
\begin{merkkasten}[Kasten 1: E-Mail an einen Austauschschüler]
\vspace{0.2em}

\textbf{Struktur einer E-Mail:}

\begin{tabularx}{\textwidth}{@{}|l|X|@{}}
\hline
\rowcolor{hellgrau}
\textbf{Teil} & \textbf{Beispiel} \\
\hline
Anrede & \es{¡Hola María! / Querido Pablo, / Hola,} \\
\hline
Vorstellung & \es{Me llamo... y tengo... años.} \\
\hline
Beschreibung & \es{Vivo en... En mi barrio hay...} \\
\hline
Vorschläge & \es{Cuando vienes, podemos ir a...} \\
\hline
Fragen & \es{¿Y tú? ¿Qué te gusta hacer?} \\
\hline
Abschluss & \es{¡Hasta pronto! / Un saludo,} \\
\hline
\end{tabularx}

\vspace{0.3em}
\textbf{Hilfreiche Strukturen:} \es{ser/estar/hay} + Orte + \es{ir a} + Infinitiv
\vspace{0.2em}
\end{merkkasten}

\vspace{1em}

% ============================================================================
% AUFGABEN
% ============================================================================

\aufgabe{Punto final: Schreibe eine E-Mail an einen Austauschschüler. (Buch S. 33)}{10}

\begin{buchverweis}
\textbf{Schreibe 8--10 Sätze.} Stell dich vor, beschreibe dein Stadtviertel und schlage Aktivitäten vor.
\end{buchverweis}

\vspace{0.3em}
\begin{hilfskasten}
\textbf{Satzanfänge:} \es{Me llamo... / Vivo en... / En mi barrio hay... / Mi barrio es... porque... / Podemos ir a... / ¿Qué te gusta...?}
\end{hilfskasten}

\vspace{0.3em}
\begin{tabularx}{\textwidth}{@{}|X|@{}}
\hline
\\[0.4em]
\luecke{14cm} \\[0.7em]
\luecke{14cm} \\[0.7em]
\luecke{14cm} \\[0.7em]
\luecke{14cm} \\[0.7em]
\luecke{14cm} \\[0.7em]
\luecke{14cm} \\[0.7em]
\luecke{14cm} \\[0.7em]
\luecke{14cm} \\[0.4em]
\hline
\end{tabularx}

\vspace{1em}

\aufgabe{Evaluación: Selbsttest -- ser/estar/hay. (Buch S. 34, Aufg. 1)}{6}

\begin{buchverweis}
\textbf{Ergänze mit der richtigen Form.}
\end{buchverweis}

\begin{enumerate}[label=\alph*), leftmargin=2em]
    \item Mi barrio \luecke{2cm} muy bonito.
    \item En el centro \luecke{2cm} muchas tiendas.
    \item La piscina \luecke{2cm} cerca de mi casa.
    \item Nosotros \luecke{2cm} de Madrid.
    \item ¿Dónde \luecke{2cm} el cine?
    \item En mi calle no \luecke{2cm} parques.
\end{enumerate}

\vspace{0.8em}

\aufgabe{Evaluación: Orte und Aktivitäten. (Buch S. 34, Aufg. 2)}{5}

\begin{buchverweis}
\textbf{Verbinde die Orte mit passenden Aktivitäten.}
\end{buchverweis}

\begin{tabularx}{\textwidth}{@{}|X|X|@{}}
\hline
\rowcolor{hellgrau}
\textbf{Ort} & \textbf{Aktivität} \\
\hline
1. la piscina & a) ver una película \\
\hline
2. el cine & b) nadar \\
\hline
3. el polideportivo & c) comprar ropa \\
\hline
4. el centro comercial & d) jugar al fútbol \\
\hline
5. la biblioteca & e) leer libros \\
\hline
\end{tabularx}

\vspace{0.3em}
Lösungen: 1 = \luecke{0.5cm}, 2 = \luecke{0.5cm}, 3 = \luecke{0.5cm}, 4 = \luecke{0.5cm}, 5 = \luecke{0.5cm}

\vspace{0.8em}

\aufgabe{Evaluación: Konjugation von ir. (Buch S. 34, Aufg. 3)}{6}

\begin{enumerate}[label=\alph*), leftmargin=2em]
    \item Yo \luecke{2.5cm} (ir) al cine los sábados.
    \item ¿Adónde \luecke{2.5cm} (ir/tú) después de clase?
    \item Nosotros \luecke{2.5cm} (ir) a la discoteca.
    \item Mis amigos \luecke{2.5cm} (ir) al parque.
    \item María \luecke{2.5cm} (ir) a la biblioteca.
    \item ¿\luecke{2.5cm} (ir/vosotros) al polideportivo?
\end{enumerate}

\vspace{0.8em}

\aufgabe{Vorbereitung Assessment: Formuliere 4 Fragen für das Partnergespräch.}{8}

\begin{buchverweis}
\textbf{Thema: Mi barrio.} Nutze: \es{¿Hay...? ¿Cómo es...? ¿Adónde vas...? ¿Qué haces...?}
\end{buchverweis}

\begin{enumerate}[leftmargin=2em]
    \item \luecke{12cm}
    \item \luecke{12cm}
    \item \luecke{12cm}
    \item \luecke{12cm}
\end{enumerate}

\vspace{1em}
\hrule
\vspace{0.3em}
\begin{flushright}
\textbf{Gesamt: \luecke{1cm} / 35 Punkte}
\end{flushright}

\end{document}
